\begin{tabular}{@{}p{0.17\linewidth}p{0.79\linewidth}@{}}
\toprule
项目 & 设置 \\
\midrule
卫星数据 & 美国国家海洋和大气管理局海岸观测网多星融合叶绿素 a 产品（经验正交函数插补去云），2 km、逐日，2018-01-01 至 2026-10-23 共 3197 天；喷气推进实验室多尺度融合海表温度，0.01° \\
任务窗口 & 每年 4--6 月浙江外海陆架（26.5°--31.5°N，121.5°--125°E）局部峰值超过 5 mg/m$^3$ 处，80 km $\times$ 80 km（40 $\times$ 40 个 2 km 网格）；开发 274 个（2018--2022 年），测试 148 个（2023--2026 年） \\
真实海洋 & 任务当天卫星图叠加卫星分辨不出的小尺度变化（标准差 0.15，相关尺度 6 km，主设置）；或任务当天卫星图本身 \\
航行器 & 航速 1.5 m/s，电池 1 kWh；平台 15 W，推进 30 W，高功耗芯片 +20 W（省电加速器 +5 W）；任务 24 h，耗尽电量即停 \\
决策周期 & 20 min；走一段约 1.8 km、采样约 2 次；停下漂流采样 5 次；荧光计单次误差 0.1 \\
想象海洋 & 48 个（32 个用于搜索航线，16 个留出），由 3 天前卫星图与数据标定的误差生成，空间相关长度 10 km，随每个样本同化更新 \\
规划器 & 交叉熵搜索经两个航路点的折线航线，跨周期热启动；低功耗处理器每分钟 5、20 或 100 次推演（条件扫描 5--200 次），高功耗芯片为其 50 倍 \\
出发前预演 & 4 个留出海洋作假想真值 $\times$ 5 种用时策略，各完整模拟 24 h，选平均假想遗憾最小者 \\
评价 & 时间平均遗憾（对数浓度，越低越好）；配对差值的 95\% 置信区间由按赤潮周整组的自助重抽样得到 \\
\bottomrule
\end{tabular}
